\documentclass[12pt]{article}
\usepackage[utf8]{inputenc}
\usepackage[T1]{fontenc}
\usepackage{hyperref}
\title{Matambwe Language Resources for K-12 Curriculum Development}
\author{Compiled by Claude Sonnet 5.5 (claude-sonnet-5-5) for Andr\'as Kornai}
\date{}
\begin{document}
\maketitle

\section{Introduction}
Matambwe (ISO 639-3 wtb) is a Bantu language of the Tunduru District, Ruvuma Region, southern Tanzania. Glottolog places it in the Makonde branch of Rufiji-Ruvuma Bantu. Wikipedia's Makonde article notes that Matambwe and Mabiha are divergent and may not be Makonde. Joshua Project gives a population of roughly 64,000 to 66,000 and says many people identify as Roman Catholic. I found no speaker count for the language itself.

\section{Primary Educational Resources}
I found no primers, textbooks or school programmes for Matambwe.

\section{Linguistic and Reference Materials}
\begin{itemize}
\item Glottolog entry: \url{https://glottolog.org/resource/languoid/id/mata1313.bigmap.html}
\item Wikipedia, Makonde language (mentions Matambwe): \url{https://en.wikipedia.org/wiki/Makonde_language}
\end{itemize}
I found no dictionary or grammar.

\section{Religious and Community Resources}
\begin{itemize}
\item Joshua Project people-group profile: \url{https://joshuaproject.net/people_groups/22372}
\end{itemize}
I found no scripture translation.

\section{Digital and Audio Resources}
I found none specific to Matambwe.

\section{Curriculum Development Context}
Documentation is thin in what I found. Tanzanian schools teach in Swahili and English. Closely related Makonde materials might help, but the relationship is disputed, so they should be used with care. Any curriculum would need new documentation work with the community.
\end{document}